\begin{titlepage}
\thispagestyle{empty}
\begin{center}
\vspace*{0.8cm}
{\Large\bfseries TRƯỜNG ĐẠI HỌC BÁCH KHOA — ĐHQG-HCM}\\[0.3cm]
{\large\bfseries KHOA KHOA HỌC VÀ KỸ THUẬT MÁY TÍNH}\\[1.2cm]

\begin{figure}[H]
\centering
% Placeholder logo trường nếu có nhu cầu chèn logo
\vspace{1.5cm}
\end{figure}

{\LARGE\bfseries BÁO CÁO BÀI TẬP LỚN}\\[0.4cm]
{\Large\bfseries MÔN HỌC: HỌC SÂU (DEEP LEARNING)}\\[1.0cm]

{\LARGE\bfseries ĐỀ TÀI (BTL03):}\\[0.3cm]
{\huge\bfseries TÁI CẤU TRÚC VÀ CHUYỂN ĐỔI MÃ NGUỒN\\ TỪ TENSORFLOW/KERAS SANG PYTORCH}\\[0.5cm]
{\Large\bfseries ĐỐI CHIẾU THỰC NGHIỆM VÀ PHÂN TÍCH TƯƠNG ĐƯƠNG MÔ HÌNH}\\[1.5cm]

\vfill

\noindent
\begin{tabularx}{\textwidth}{c >{\raggedright\arraybackslash\ttfamily\small}p{0.37\textwidth} >{\centering\arraybackslash}p{0.20\textwidth} >{\raggedright\arraybackslash}X}
\toprule
\textbf{STT} & \textrm{\textbf{Mã thành viên}} & \textbf{Họ tên / MSSV} & \textbf{Phân công nhiệm vụ} \\
\midrule
1 & Duy-DS & (Chờ cập nhật) & Trưởng nhóm; Hạ tầng OOP \& Mô hình MLP (WP01, WP02) \\
2 & nguyenthanhdat180905-cloud & (Chờ cập nhật) & Tiền xử lý dữ liệu NCHW \& Mô hình CNN+MLP (WP02) \\
3 & nguyen-cong-tri & (Chờ cập nhật) & Dữ liệu chuỗi ảnh \& Mô hình CNN+LSTM (WP02, WP03) \\
4 & TonOfWeight072501 & (Chờ cập nhật) & Hệ thống Callbacks, Logging \& Ablation Studies (WP02, WP03) \\
5 & CHIVY2005 & (Chờ cập nhật) & Mô hình Thị giác nâng cao SegNet/FCN \& Báo cáo (WP02, WP04) \\
\midrule
\multicolumn{2}{@{}l@{}}{\small\textbf{Giảng viên hướng dẫn:} (Chờ cập nhật)} & \multicolumn{2}{r@{}}{\small\textbf{Đóng góp toàn nhóm:} 100\% (đồng đều)} \\
\bottomrule
\end{tabularx}

\vfill
{\large Thành phố Hồ Chí Minh, Năm 2026}
\end{center}
\end{titlepage}

\chapter*{Lời cam đoan và Lời cảm ơn}
\addcontentsline{toc}{chapter}{Lời cam đoan và Lời cảm ơn}
Chúng em xin cam đoan rằng toàn bộ nội dung trong báo cáo này là kết quả nghiên cứu, phân tích mã nguồn và thực nghiệm thực tế do chính tập thể nhóm thực hiện dưới sự hướng dẫn của giảng viên phụ trách môn học. Mọi số liệu thực nghiệm, kết quả đối chiếu số lượng tham số giữa Keras và PyTorch, cùng các trích dẫn tài liệu tham khảo đều được ghi nhận trung thực dựa trên mã nguồn và nhật ký thực nghiệm lưu trữ trong kho lưu trữ của nhóm.

Nhóm xin chân thành cảm ơn Thầy/Cô bộ môn Học sâu đã tạo điều kiện và hướng dẫn tận tình, giúp chúng em có cơ hội đào sâu bản chất toán học cũng như cơ chế thực thi bên dưới của hai framework học sâu phổ biến hàng đầu hiện nay.

\chapter*{Tóm tắt báo cáo}
\addcontentsline{toc}{chapter}{Tóm tắt báo cáo}
Đề tài Bài tập lớn môn Học sâu (BTL03) hướng tới mục tiêu nghiên cứu, tái cấu trúc và chuyển đổi hoàn chỉnh 10 mô hình học sâu thuộc 5 nhóm kiến trúc kinh điển từ nền tảng TensorFlow/Keras sang thư viện PyTorch thuần túy theo chuẩn lập trình hướng đối tượng (OOP). Các nhóm bài toán bao gồm: Mạng nơ-ron truyền thẳng (MLP), Mạng tích chập kết hợp MLP (CNN+MLP), Mạng kết hợp chuỗi thời gian (CNN+LSTM), Hệ thống tiện ích huấn luyện (Callbacks/Logger), và Thị giác máy tính nâng cao (Phân vùng ảnh SegNet và Phát hiện vật thể FCN).

Trong quá trình thực hiện, nhóm tập trung giải quyết các nút thắt kỹ thuật trọng tâm: chuyển đổi trật tự chiều tensor ảnh từ NHWC sang NCHW, bảo toàn phép tiền xử lý xoay ảnh không gian, thay thế lớp \texttt{TimeDistributed} của Keras bằng cơ chế gộp chiều tensor qua CNN và phân tách lại cho LSTM, đồng thời tái lập chuẩn xác cấu trúc bộ giải mã phân vùng ảnh SegNet. Về mặt thực nghiệm, nhóm đã kiểm chứng tính tương đương chính xác về số lượng tham số (Parameter Parity) trên 10/10 mô hình, thực hiện khảo sát tham số sweep trên mô hình xấp xỉ hàm (T0201) và huấn luyện 30 epoch trên tập MNIST (T0202). Trong đợt kiểm thử tích hợp tự động toàn hệ thống (T0212), cả 10 mô hình đều vượt qua thành công chu trình lan truyền xuôi và ngược trên dữ liệu kiểm thử nhanh; 78/82 bài kiểm tra chức năng đạt yêu cầu, và nhóm đang tiếp tục tinh chỉnh môi trường logging trước khi tiến hành thực nghiệm chuyên sâu trên toàn bộ tập dữ liệu thực tế.

\noindent\textbf{Từ khóa:} Học sâu; TensorFlow/Keras; PyTorch; Chuyển đổi kiến trúc; Định dạng NCHW; Lập trình hướng đối tượng; Tái lập thực nghiệm.

\chapter*{Danh mục từ viết tắt}
\addcontentsline{toc}{chapter}{Danh mục từ viết tắt}
\begin{tabularx}{\textwidth}{lY}
\toprule \textbf{Viết tắt} & \textbf{Thuật ngữ tiếng Anh và Giải thích ý nghĩa} \\
\midrule
CNN & Convolutional Neural Network — Mạng nơ-ron tích chập trích xuất đặc trưng không gian \\
FCN & Fully Convolutional Network — Mạng tích chập toàn phần ứng dụng hồi quy bản đồ điểm \\
LSTM & Long Short-Term Memory — Mạng bộ nhớ ngắn hạn dài xử lý dữ liệu chuỗi \\
MLP & Multi-Layer Perceptron — Mạng nơ-ron truyền thẳng đa tầng \\
NCHW & Batch -- Channel -- Height -- Width (Định dạng tensor ảnh chuẩn trong PyTorch) \\
NHWC & Batch -- Height -- Width -- Channel (Định dạng tensor ảnh mặc định trong Keras) \\
OOP & Object-Oriented Programming — Phương pháp lập trình hướng đối tượng \\
RQ & Research Question — Câu hỏi nghiên cứu thực nghiệm của đề tài \\
WBS & Work Breakdown Structure — Cấu trúc phân chia công việc dự án \\
\bottomrule
\end{tabularx}

\chapter*{Câu hỏi nghiên cứu thực nghiệm}
\addcontentsline{toc}{chapter}{Câu hỏi nghiên cứu thực nghiệm}
Nhằm đánh giá sâu sắc sự khác biệt bản chất giữa TensorFlow/Keras và PyTorch, đề tài đặt ra hai câu hỏi nghiên cứu trọng tâm làm kim chỉ nam cho quá trình thiết kế thực nghiệm:

\begin{longtable}{L{0.14\textwidth}L{0.38\textwidth}L{0.28\textwidth}L{0.16\textwidth}}
\toprule \textbf{Mã RQ} & \textbf{Nội dung câu hỏi nghiên cứu} & \textbf{Phương pháp kiểm chứng} & \textbf{Hiện trạng} \\
\midrule
RQ1 & Sự khác biệt về định dạng tensor (NHWC vs. NCHW), vòng lặp huấn luyện tường minh và cơ chế khởi tạo trọng số mặc định ảnh hưởng ra sao đến tốc độ hội tụ và độ ổn định gradient? & Thiết lập thực nghiệm đối chứng trên cùng bộ tham số siêu dữ liệu, hạt giống ngẫu nhiên và bộ tối ưu hóa. & Đã hoàn thành bước chuyển đổi kiến trúc; đang chuẩn bị thực nghiệm đối chứng. \\
RQ2 & Các kỹ thuật điều chuẩn (Dropout, BatchNorm) và tiện ích huấn luyện (Data Augmentation, EarlyStopping, ReduceLR) kiểm soát hiện tượng quá khớp (overfitting) hiệu quả như thế nào trên dữ liệu nhiều nhiễu? & Tiến hành nghiên cứu cắt bỏ (Ablation Studies ABL-01 đến ABL-03) trên tập dữ liệu nhận diện cảm xúc khuôn mặt (FER). & Đã xây dựng đầy đủ mã nguồn tiện ích; chờ kết quả huấn luyện dữ liệu thực. \\
\bottomrule
\end{longtable}

\chapter*{Lời mở đầu}
\addcontentsline{toc}{chapter}{Lời mở đầu}
\section*{Bối cảnh và Động lực đề tài}
Trong những năm gần đây, học sâu (Deep Learning) đã phát triển vượt bậc và trở thành công nghệ cốt lõi trong nhiều bài toán thực tế. Hai framework được sử dụng phổ biến nhất trong nghiên cứu và công nghiệp là TensorFlow/Keras và PyTorch. Trong khi Keras nổi bật nhờ khả năng xây dựng mô hình nhanh chóng với mức độ trừu tượng hóa cao, thì PyTorch lại chiếm ưu thế mạnh mẽ trong cộng đồng nghiên cứu khoa học nhờ đồ thị tính toán động, vòng lặp huấn luyện tường minh và sự linh hoạt tối đa trong việc kiểm soát tensor.

Việc chuyển đổi các kiến trúc mô hình từ Keras sang PyTorch không chỉ đơn thuần là dịch cú pháp lệnh, mà còn đòi hỏi người kỹ sư phải hiểu thấu đáo các khác biệt về mặt cấu trúc toán học: sự hoán vị trục không gian của dữ liệu ảnh, sự khác biệt trong cách tính toán hàm mất mát phân loại, cơ chế lan truyền ngược và cách thức quản lý trạng thái mô hình. Đây chính là động lực để nhóm chúng em thực hiện đề tài Bài tập lớn 03.

\section*{Mục tiêu và Phạm vi nghiên cứu}
Mục tiêu bao trùm của đề tài là tái hiện chính xác 10 mô hình học sâu thuộc 5 nhóm bài toán từ Keras sang 100\% mã nguồn PyTorch hướng đối tượng, đảm bảo tính tương đương toán học về số lượng tham số huấn luyện và khả năng tái lập thực nghiệm. Báo cáo này tổng kết toàn bộ quá trình nghiên cứu, giải thích cặn kẽ các giải pháp kỹ thuật, trình bày các kết quả kiểm thử chức năng đã đạt được và nêu rõ kế hoạch hoàn thiện đề tài.

\section*{Bố cục báo cáo}
Nội dung báo cáo được tổ chức thành 5 chương chính:
\begin{itemize}
\item \textbf{Chương 1 — Tổng quan bài toán và phạm vi chuyển đổi:} Giới thiệu 10 mô hình thực hành, phân chia công việc theo WBS và hiện trạng tiến độ.
\item \textbf{Chương 2 — Cơ sở lý thuyết và nguyên lý chuyển đổi Keras sang PyTorch:} Phân tích cơ chế toán học và giải pháp xử lý các điểm khác biệt cốt lõi giữa hai framework.
\item \textbf{Chương 3 — Thiết kế kiến trúc hệ thống và quy trình chuyển đổi:} Trình bày thiết kế hướng đối tượng, sơ đồ luồng dữ liệu, hệ thống callbacks và giao diện thực thi CLI.
\item \textbf{Chương 4 — Kết quả chuyển đổi, kiểm thử tích hợp và hiện trạng thực nghiệm:} Trình bày bảng đối chiếu tham số, kết quả sweep/huấn luyện ban đầu và phân tích kiểm thử.
\item \textbf{Chương 5 — Phân tích độ tin cậy, bài học kinh nghiệm và hướng phát triển:} Đánh giá những ưu nhược điểm, rút ra bài học thực tiễn và vạch ra lộ trình hoàn thiện cho giai đoạn kế tiếp.
\end{itemize}
